%% ============================================================
%% RESUMES TRI-LANGUE -- une seule page, derniere du rapport
%% FR / EN / AR -- chaque section <= 4 lignes + mots-cles
%% ============================================================
\cleardoublepage
\thispagestyle{empty}
\pagestyle{empty}

\begingroup
\setlength{\parskip}{0pt}
\setlength{\parindent}{0pt}
\renewcommand{\baselinestretch}{1.1}
\footnotesize

% ---------- Français ----------
\begin{center}
  \textbf{\normalsize Résumé}
\end{center}
\noindent\textbf{Mots-clés :} IoT, sécurité, TLS~1.3, mTLS, PKI, HashiCorp Vault, MQTT, Suricata, Wazuh, ML, GNS3.

\smallskip
\noindent Ce projet, mené avec \TT{}, conçoit et valide une architecture de sécurisation de bout en bout d'un
écosystème IoT simulé sous GNS3. La solution associe une PKI HashiCorp Vault à deux niveaux, un broker MQTT
en TLS~1.3/mTLS obligatoire, un IDS Suricata et un SIEM Wazuh. Un pipeline ML (IsolationForest + RandomForest)
détecte les anomalies : F1\,=\,0,994 et ROC-AUC\,=\,0,999 pour RF ; surcoût TLS de +89\,\% sur le RTT MQTT
(4,2 $\to$ 7,9\,ms), en deçà du seuil critique de 50\,ms admis pour l'IoT industriel.

\vspace{4pt}\noindent\rule{\linewidth}{0.4pt}\vspace{4pt}

% ---------- English ----------
\begin{center}
  \textbf{\normalsize Abstract}
\end{center}
\noindent\textbf{Keywords:} IoT, security, TLS~1.3, mTLS, PKI, HashiCorp Vault, MQTT, Suricata, Wazuh, ML, GNS3.

\smallskip
\noindent This project, carried out with Tunisie Telecom, designs and validates an end-to-end security architecture for a
simulated IoT ecosystem under GNS3. The solution combines a two-level HashiCorp Vault PKI, a Mosquitto MQTT
broker enforcing TLS~1.3/mTLS, a Suricata IDS and a Wazuh SIEM. An ML pipeline (IsolationForest + RandomForest)
performs anomaly detection: F1\,=\,0.994, ROC-AUC\,=\,0.999 (RF); TLS overhead is +89\,\% on RTT
(4.2 $\to$ 7.9\,ms), well below the 50\,ms threshold for industrial IoT.

\vspace{4pt}\noindent\rule{\linewidth}{0.4pt}\vspace{4pt}

% ---------- Arabic ----------
\begin{center}
  \textbf{\normalsize \textarabic{ملخّص}}
\end{center}
\begin{Arabic}
\noindent\textbf{الكلمات المفتاحيّة:}
إنترنت الأشياء، أمن سيبراني، TLS~1.3، mTLS، PKI، HashiCorp Vault، MQTT، Suricata، Wazuh، تعلّم آليّ، GNS3.

\smallskip
\noindent صمّم هذا المشروع، المُنجَز مع \emph{اتّصالات تونس}، بنيةً لتأمين الاتّصالات من طرف إلى طرف داخل منظومة
إنترنت الأشياء المُحاكاة عبر GNS3. تجمع الحلّ بين بنية مفاتيح عموميّة (HashiCorp Vault)، ووسيط MQTT بـ TLS~1.3/mTLS،
ونظام Suricata، ومنصّة Wazuh SIEM. يُحقّق نموذج RandomForest في سلسلة التعلّم الآليّ
F1\,=\,0,994 وROC-AUC\,=\,0,999؛ وزيادة TLS على زمن الرحلة الذهابيّة-الإيابيّة +89\,\%
(4,2 $\to$ 7,9 ميلّي ثانية) — أدنى من العتبة المقبولة (50 ميلّي ثانية) للإنترنت الصناعيّ.
\end{Arabic}

\endgroup
%% Resume tri-langue compact -- doit tenir sur UNE seule page (derniere page du rapport)
\cleardoublepage
\thispagestyle{empty}
\pagestyle{empty}

\begingroup
\setlength{\parskip}{2pt}
\setlength{\parindent}{0pt}
\small

%% ----------- FR -----------
\noindent\textbf{\large Résumé}\\[2pt]
\textbf{Mots-clés :} IoT, sécurité, TLS 1.3, mTLS, PKI, HashiCorp Vault, MQTT, Suricata, Wazuh, SIEM, Machine Learning, GNS3, Scrum.\\[2pt]
Ce projet, réalisé en partenariat avec \TT{}, propose une architecture de sécurisation de bout en bout des flux d'un écosystème IoT simulé sous GNS3. La solution combine une PKI HashiCorp Vault à deux niveaux, un broker MQTT Mosquitto en TLS~1.3/mTLS, un simulateur de 50 dispositifs rejouant 76\,000 messages, un IDS Suricata, un SIEM Wazuh et un pipeline ML (IsolationForest + RandomForest). Les résultats : F1=0,994 / ROC-AUC=0,999 (RF) ; surcoût TLS de +89\,\% sur le RTT (4,2$\to$7,9~ms), bien sous le seuil de 50~ms admis pour l'IoT industriel.

\vspace{6pt}\hrule\vspace{6pt}

%% ----------- EN -----------
\noindent\textbf{\large Abstract}\\[2pt]
\textbf{Keywords:} IoT, security, TLS 1.3, mTLS, PKI, HashiCorp Vault, MQTT, Suricata, Wazuh, SIEM, Machine Learning, GNS3, Scrum.\\[2pt]
This project, carried out with Tunisie Telecom, proposes an end-to-end security architecture for a simulated IoT ecosystem under GNS3. The solution combines a two-level HashiCorp Vault PKI, a Mosquitto MQTT broker enforcing TLS~1.3/mTLS, a 50-device fleet simulator replaying 76\,000 messages, a Suricata IDS, a Wazuh SIEM and an ML pipeline (IsolationForest + RandomForest). Results: F1=0.994 / ROC-AUC=0.999 (RF); TLS overhead of +89\,\% on RTT (4.2$\to$7.9~ms), well below the 50~ms threshold accepted for industrial IoT.

\vspace{6pt}\hrule\vspace{6pt}

%% ----------- AR -----------
\noindent\textbf{\large \textarabic{ملخّص}}\\[2pt]
\begin{Arabic}
\textbf{الكلمات المفتاحيّة:} إنترنت الأشياء، الأمن السيبراني، TLS~1.3، mTLS، PKI، HashiCorp Vault، MQTT، Suricata، Wazuh، التعلّم الآليّ، GNS3، Scrum.\\[2pt]
يقترح هذا المشروع، المُنجَز بالشراكة مع \emph{اتّصالات تونس}، بنيةً لتأمين الاتّصالات من طرفٍ إلى طرف داخل منظومة إنترنت الأشياء مُحاكاة عبر GNS3. تجمع الحلّ بين بنية مفاتيح عموميّة على HashiCorp Vault بمستويَين، ووسيط MQTT (Mosquitto) بـ TLS~1.3/mTLS، ومحاكي 50 جهازًا يُعيد 76\,000 رسالة، ونظام Suricata، وSIEM Wazuh، وسلسلة تعلّم آليّ (IsolationForest + RandomForest). النتائج: F1=0,994 وROC-AUC=0,999 (RF)؛ زيادة TLS بنسبة 89\,\% في زمن الذهاب والإياب (من 4,2 إلى 7,9 ميلّي ثانية)، أي أدنى بكثير من العتبة 50 ميلّي ثانية المقبولة في إنترنت الأشياء الصناعيّ.
\end{Arabic}

\endgroup
\cleardoublepage
\chapter*{Résumé}
\addcontentsline{toc}{chapter}{Résumés}
\markboth{Résumés}{Résumés}

\noindent\textbf{Mots-clés :} IoT, sécurité, TLS 1.3, mTLS, PKI, HashiCorp Vault, MQTT, Suricata, Wazuh, SIEM, IsolationForest, RandomForest, détection d'anomalies, GNS3, Scrum.

\noindent L'essor de l'Internet des Objets (IoT) dans les environnements des opérateurs de télécommunications engendre des flux de données massifs et hétérogènes exposés à des risques croissants d'interception, d'altération et de compromission. Ce projet de fin d'études, réalisé en partenariat avec \TT{}, aborde la problématique de la \textbf{sécurisation de bout en bout des flux de communication d'un écosystème IoT simulé}.

\noindent La démarche adoptée suit la méthodologie \textbf{Scrum} et s'appuie sur un laboratoire simulé combinant \textbf{GNS3} pour la topologie réseau et \textbf{Docker} pour les nœuds applicatifs. Nous avons conçu et déployé une architecture composée de six éléments majeurs : (1)~une \textbf{infrastructure à clés publiques (PKI)} basée sur HashiCorp Vault avec une hiérarchie CA racine -- CA intermédiaire ; (2)~un \textbf{broker MQTT Mosquitto} configuré en TLS~1.3 et mTLS obligatoire ; (3)~un \textbf{simulateur de flotte IoT} rejouant 76\,000 événements réels couvrant sept types d'attaques ; (4)~une \textbf{analyse réseau} par Suricata avec règles personnalisées ; (5)~un \textbf{SIEM Wazuh} ingérant les alertes et corrélant les événements ; (6)~un \textbf{pipeline de détection d'anomalies} par Machine Learning (IsolationForest + RandomForest).

\noindent Les résultats obtenus démontrent l'efficacité de l'approche : le RandomForest atteint un F1-score de \textbf{0,994} et un ROC-AUC de \textbf{0,999}, tandis que l'IsolationForest, en mode non supervisé, obtient \textbf{0,930} / \textbf{0,959}. L'overhead TLS~1.3/mTLS sur le RTT MQTT est de \textbf{+89\,\%} (4,2\,ms $\to$ 7,9\,ms), restant largement en deçà du seuil critique de 50\,ms pour les applications IoT industrielles.

\bigskip

\chapter*{Abstract}
\addcontentsline{toc}{chapter}{Abstract}

\noindent\textbf{Keywords:} IoT, security, TLS 1.3, mTLS, PKI, HashiCorp Vault, MQTT, Suricata, Wazuh, SIEM, IsolationForest, RandomForest, anomaly detection, GNS3, Scrum.

\noindent The growth of the Internet of Things (IoT) in telecom operator environments generates massive and heterogeneous data flows exposed to increasing risks of interception, alteration and compromise. This final-year project, carried out in partnership with Tunisie Telecom and following the Scrum methodology, addresses the \textbf{end-to-end security of communication flows in a simulated IoT ecosystem}.

\noindent We designed and deployed a complete architecture combining \textbf{GNS3} for network topology and \textbf{Docker} for application nodes, and including six major components: a HashiCorp Vault PKI with a two-level CA hierarchy, a Mosquitto MQTT broker enforcing TLS~1.3 and mTLS, a 50-device IoT fleet simulator replaying 76\,000 real events across seven attack categories, a Suricata IDS with custom MQTT rules, a Wazuh SIEM correlating the alerts, and a Machine Learning anomaly detection pipeline based on IsolationForest and RandomForest.

\noindent The RandomForest classifier achieves an F1-score of \textbf{0.994} and a ROC-AUC of \textbf{0.999} on binary classification, while the unsupervised IsolationForest reaches \textbf{0.930} / \textbf{0.959}. TLS overhead on the MQTT round-trip time is \textbf{+89\,\%} (4.2\,ms $\to$ 7.9\,ms), which remains well below the 50\,ms threshold accepted for industrial IoT applications.

\bigskip

\chapter*{\textarabic{ملخّص}}
\addcontentsline{toc}{chapter}{\textarabic{ملخّص}}
\markboth{Resume arabe}{Resume arabe}

\begin{Arabic}
\noindent\textbf{الكلمات المفتاحيّة:} إنترنت الأشياء، الأمن السيبراني، بروتوكول TLS~1.3، المصادقة المتبادلة mTLS، البنية التحتيّة للمفاتيح العموميّة PKI، HashiCorp Vault، MQTT، Suricata، Wazuh، التعلّم الآليّ، كشف الشذوذ، GNS3، منهجيّة Scrum.

\noindent يشهد انتشار إنترنت الأشياء (IoT) لدى مشغّلي الاتّصالات تدفّقًا هائلًا ومتنوّعًا للبيانات، ممّا يجعلها عرضةً لمخاطر متزايدة من الاعتراض والتعديل والاختراق. يتناول هذا المشروع، المُنجَز بالشراكة مع \emph{اتّصالات تونس}، إشكاليّةَ \textbf{تأمين الاتّصالات من طرفٍ إلى طرف داخل منظومة إنترنت الأشياء} في بيئة محاكاة.

\noindent اعتمدنا منهجيّة \textbf{Scrum} في إدارة المشروع، وأنشأنا مختبرًا افتراضيًّا يجمع بين GNS3 لمحاكاة الشبكة وDocker لحاويات الخدمات. تتألّف البنية المُقترَحة من ستّة مكوّنات رئيسيّة: بنية تحتيّة للمفاتيح العموميّة مبنيّة على HashiCorp Vault بتسلسل هرميّ من شهادتَي CA؛ ووسيط MQTT (Mosquitto) مُهيَّأ بـ TLS~1.3 ومصادقة متبادلة إجباريّة؛ ومحاكي أسطول مكوّن من 50 جهازًا يُعيد تشغيل 76\,000 حدث حقيقيّ يغطّي سبعة أنواع من الهجمات؛ ونظام كشف تسلّل Suricata بقواعد مخصّصة لبروتوكول MQTT؛ ومنصّة SIEM (Wazuh) لتجميع التنبيهات والربط بينها؛ وأخيرًا سلسلة لكشف الشذوذ بالتعلّم الآليّ تجمع بين خوارزميتَي IsolationForest وRandomForest.

\noindent أثبتت النتائج فعاليّة المقاربة: حقّقت خوارزميّة RandomForest قيمة F1 تساوي \textbf{0,994} وقيمة ROC-AUC تساوي \textbf{0,999} في التصنيف الثنائي، في حين بلغت IsolationForest غير المُشرَف عليها \textbf{0,930} و\textbf{0,959}. أمّا الزيادة في زمن التأخّر الناجمة عن TLS~1.3/mTLS فبلغت \textbf{89\,\%} (من 4{,}2 إلى 7{,}9 ميلّي ثانية)، وهي قيمة تظلّ أدنى بكثير من العتبة الحرجة (50 ميلّي ثانية) المقبولة في تطبيقات إنترنت الأشياء الصناعيّة.
\end{Arabic}

\cleardoublepage
\chapter*{Résumé}
\addcontentsline{toc}{chapter}{Résumé}

\textbf{Mots-clés :} IoT, sécurité, TLS 1.3, mTLS, PKI, HashiCorp Vault, MQTT, Suricata, Wazuh, SIEM, IsolationForest, RandomForest, détection d'anomalies, GNS3, Scrum.

\vspace{0.5cm}

L'essor de l'Internet des Objets (IoT) dans les environnements des opérateurs de télécommunications engendre des flux de données massifs et hétérogènes exposés à des risques croissants d'interception, d'altération et de compromission. Ce projet de fin d'études, réalisé en partenariat avec \TT{}, aborde la problématique de la \textbf{sécurisation de bout en bout des flux de communication d'un écosystème IoT simulé}.

La démarche adoptée suit la méthodologie \textbf{Scrum} et s'appuie sur un laboratoire simulé combinant \textbf{GNS3} pour la topologie réseau et \textbf{Docker} pour les nœuds applicatifs. Nous avons conçu et déployé une architecture composée de six éléments majeurs : (1)~une \textbf{infrastructure à clés publiques (PKI)} basée sur HashiCorp Vault avec une hiérarchie CA racine -- CA intermédiaire ; (2)~un \textbf{broker MQTT Mosquitto} configuré en TLS~1.3 et mTLS obligatoire ; (3)~un \textbf{simulateur de flotte IoT} rejouant 76\,000 événements réels issus d'un dataset CSV couvrant sept types d'attaques ; (4)~une \textbf{analyse réseau} par Suricata avec règles personnalisées ; (5)~un \textbf{SIEM Wazuh} ingérant les alertes et corrélant les événements ; (6)~un \textbf{pipeline de détection d'anomalies} par Machine Learning (IsolationForest + RandomForest).

Les résultats obtenus démontrent l'efficacité de l'approche : le RandomForest atteint un F1-score de \textbf{0,994} et un ROC-AUC de \textbf{0,999} pour la détection binaire, tandis que l'IsolationForest, en mode non supervisé, obtient \textbf{0,930} / \textbf{0,959}. Les tests de performance montrent que l'overhead TLS~1.3/mTLS sur le RTT MQTT est de \textbf{+89\,\%} (4,2\,ms $\to$ 7,9\,ms), restant largement en deçà du seuil critique de 50\,ms pour les applications IoT industrielles.

\vspace{0.5cm}
\hrule
\vspace{0.3cm}

\textbf{Abstract} -- \textit{The growth of IoT in telecom operator environments generates massive and heterogeneous data flows exposed to increasing risks of interception, alteration and compromise. This final-year project, carried out in partnership with Tunisie Telecom and following Scrum methodology, addresses the end-to-end security of communication flows in a simulated IoT ecosystem. We designed and deployed a complete architecture including a PKI based on HashiCorp Vault, a Mosquitto MQTT broker with TLS~1.3/mTLS, an IoT fleet simulator, Suricata IDS, Wazuh SIEM and an ML anomaly detection pipeline. The RandomForest classifier achieves F1=0.994, ROC-AUC=0.999. TLS overhead on RTT is +89\,\% (4.2\,ms to 7.9\,ms), well within IoT industrial thresholds.}

\cleardoublepage
